%==============================================================================
%  CHAPITRE 5 — VALIDATION ET RÉSULTATS
%==============================================================================
\titreduchapitre{Validation et Résultats}{chap:validation}

Ce chapitre présente l'évaluation multi-niveaux du prototype \emolearn{},
conformément à l'OS6. Il distingue les validations \textbf{techniques}
(conduites et chiffrées) des validations \textbf{expertes et utilisateurs}
(instrumentées mais non finalisées dans le cadre de ce mémoire).

%==============================================================================
\section{Méthodologie d'évaluation multi-niveaux}
\label{sec:val-methodo}

\subsection{Approche globale}

L'évaluation s'organise en quatre niveaux complémentaires :

\begin{enumerate}[leftmargin=1.4em]
  \item \textbf{Niveau technique — perception.} Exactitude des modules visage
  (FER\,$\rightarrow$\,DAiSEE), voix (SER) et fusion multimodale sur jeux de
  test documentés.

  \item \textbf{Niveau décision.} Cohérence du mapping
  engagement\,$\rightarrow$\,stratégie et couverture des cinquante scénarios
  simulés.

  \item \textbf{Niveau contenu.} Métriques automatiques de forme sur le corpus
  généré ; protocole d'évaluation experte préparé (60 fiches).

  \item \textbf{Niveau utilisateur.} Protocole d'étude pilote ($n=30$) documenté,
  non exécuté dans la version actuelle du mémoire.
\end{enumerate}

Cette hiérarchie privilégie la \emph{mesure honnête} des compromis observés
plutôt qu'une démonstration de performance maximale.

\subsection{Métriques et critères}

Les métriques retenues varient selon le niveau :

\begin{itemize}[leftmargin=1.4em]
  \item \textbf{Classification} : accuracy, F1 macro, matrices de confusion ;
  pour la fusion, rappel des états « urgents » (confusion, frustration).
  \item \textbf{Corpus généré} : longueur en mots, indice de lisibilité
  (Kandel--Moles), taux de répétition, taux de conformité à l'audit automatique.
  \item \textbf{Évaluation experte} (protocole) : Likert 1--5 sur pertinence,
  clarté, adéquation émotionnelle, qualité pédagogique, utilisabilité MOOC et
  absence d'erreurs factuelles (avec RAG).
\end{itemize}

%==============================================================================
\section{Validation de la perception et de la fusion}
\label{sec:val-perception}

\subsection{Classifieur FER\,$\rightarrow$\,DAiSEE (Random Forest)}

Un échantillon de \textbf{1\,200} clips \textsc{DAiSEE} a servi à l'entraînement
d'un Random Forest (GridSearchCV, $n\_estimators=200$) sur les sept scores de
probabilité du modèle FER pré-entraîné \cite{li2020deepfer,wolf2020transformers}. Le jeu de test ($n=300$) est fortement
déséquilibré ($\approx$75\,\% Engagement).

\begin{table}[H]
\centering
\small
\renewcommand{\arraystretch}{1.2}
\begin{tabular}{|l|c|c|}
\hline
\rowcolor{gray!20}
\textbf{Approche} & \textbf{Accuracy} & \textbf{F1 macro} \\
\hline
Mapping heuristique (baseline) & 0{,}193 & 0{,}141 \\
\hline
\textbf{Random Forest appris} & \textbf{0{,}723} & \textbf{0{,}264} \\
\hline
Gain absolu & +0{,}530 & +0{,}123 \\
\hline
\end{tabular}
\caption{Comparaison mapping heuristique vs classifieur appris (FER\,$\rightarrow$\,DAiSEE, $n=300$).}
\label{tab:fer-rf}
\end{table}

Le classifieur appris améliore nettement l'accuracy mais le F1 macro reste
modeste : le modèle prédit surtout \emph{Engagement}, avec un rappel nul sur
\emph{Confusion} et \emph{Frustration} sur ce split. Le gain justifie néanmoins
le déploiement du modèle appris en production, sous réserve de compléter
l'évaluation par le rappel des classes minoritaires.

\subsection{Fine-tuning SER (tête wav2vec2)}

La tête de classification d'un modèle \emph{wav2vec2} pré-entraîné
\cite{baevski2020wav2vec2} a été fine-tunée sur 600 segments \textsc{IEMOCAP}
(150 par état DAiSEE, oversampling des classes rares), avec validation interne
$n=120$ sur deux époques (CPU) :

\begin{table}[H]
\centering
\small
\begin{tabular}{|c|c|c|}
\hline
\rowcolor{gray!20}
\textbf{Époque} & \textbf{Accuracy} & \textbf{F1 macro} \\
\hline
1 & 0{,}275 & 0{,}279 \\
\hline
2 & 0{,}300 & \textbf{0{,}309} \\
\hline
\end{tabular}
\caption{Progression du fine-tuning SER (validation interne).}
\label{tab:ser-ft}
\end{table}

Ces chiffres mesurent la progression d'entraînement ; la performance en usage
réel est évaluée ci-dessous sur le banc d'essai audiovisuel aligné.

\subsection{Banc d'essai audiovisuel aligné ($n=80$)}

Un sous-corpus \textsc{IEMOCAP} \textbf{aligné} (même locuteur, visage et voix
synchronisés, 20 segments par état) permet de comparer modalités uniques et
fusion sans biais d'appariement. Le tableau~\ref{tab:fusion-av} synthétise les
résultats.

\begin{table}[H]
\centering
\small
\renewcommand{\arraystretch}{1.15}
\begin{tabular}{|l|c|c|}
\hline
\rowcolor{gray!20}
\textbf{Variante} & \textbf{Accuracy} & \textbf{Rappel urgents} \\
\hline
Visage — FER 1 frame & 0{,}375 & 0{,}250 \\
\hline
Visage — FER multi-frames & 0{,}412 & 0{,}275 \\
\hline
Voix — SER heuristique & 0{,}263 & 0{,}225 \\
\hline
Voix — \textbf{SER appris (wav2vec2)} & \textbf{0{,}487} & 0{,}200 \\
\hline
Fusion — priorité pédagogique & 0{,}338 & \textbf{0{,}400} \\
\hline
Fusion — \textbf{vote pondéré} & \textbf{0{,}500} & 0{,}125 \\
\hline
\end{tabular}
\caption{Validation multimodale sur sous-corpus IEMOCAP aligné ($n=80$).}
\label{tab:fusion-av}
\end{table}

\paragraph{Lecture.} Le SER appris double quasi l'accuracy de la voix
($+0{,}224$ vs heuristique). La fusion par \emph{weighted vote} atteint la
meilleure accuracy (0{,}500) ; la stratégie \emph{pedagogical priority} maximise
le rappel des états urgents (0{,}400) au prix d'une accuracy plus faible. Ce
compromiss est assumé : le prototype retient 60\,\%/40\,\% (visage/voix) pour
équilibrer les deux objectifs (tableau~\ref{tab:poids-fusion}).

\begin{table}[H]
\centering
\small
\begin{tabular}{|c|c|c|}
\hline
\rowcolor{gray!20}
\textbf{Poids visage / voix} & \textbf{Accuracy} & \textbf{Rappel urgents} \\
\hline
0{,}5 / 0{,}5 & 0{,}287 & 0{,}375 \\
\hline
\textbf{0{,}6 / 0{,}4 (retenu)} & 0{,}300 & \textbf{0{,}425} \\
\hline
0{,}7 / 0{,}3 & 0{,}338 & 0{,}400 \\
\hline
0{,}8 / 0{,}2 & 0{,}325 & 0{,}350 \\
\hline
\end{tabular}
\caption{Balayage des poids de fusion (variante vote pondéré, $n=80$).}
\label{tab:poids-fusion}
\end{table}

\subsection{Analyse de sensibilité du mapping IEMOCAP\,$\rightarrow$\,DAiSEE}

Sur 10\,039 utterances \textsc{IEMOCAP}, des variantes plausibles du mapping
émotion\,$\rightarrow$\,engagement modifient jusqu'à \textbf{17\,\%} des
étiquettes (neutre\,$\rightarrow$\,Engagement au lieu d'Ennui). La sensibilité
confirme que cette hypothèse n'est pas anodine et justifie une validation
experte prioritaire sur les correspondances neutre/tristesse.

%==============================================================================
\section{Validation du modèle d'adaptation}
\label{sec:val-decision}

\subsection{Analyse de la pertinence du mapping}

La cascade engagement\,$\rightarrow$\,type de contenu (A/B/C/D) + patron (A1--D5)
a été appliquée aux cinquante scénarios du référentiel. La distribution des
patrons est concentrée sur l'encouragement (type D) en cas de frustration
élevée, conforme à la conception (section~\ref{sec:impl-corpus} du
chapitre~\ref{chap:implementation}).

\subsection{Couverture des états d'engagement}

Les cinquante séquences couvrent les quatre états \textsc{DAiSEE} avec un
échantillonnage stratifié favorisant les états problématiques (confusion,
frustration). Six patrons de séquence (stable, dégradation, récupération, etc.)
sont représentés.

\subsection{Validation par experts pédagogiques}

Un protocole Delphi simplifié (deux tours, trois experts, 15 cas tests +
20 fiches de patrons) a été préparé. Les fiches d'évaluation
(\texttt{results/validation\_expert\_fiches.csv}, 58 contenus stratifiés) sont
prêtes ; les scores Likert restent à compléter. Le livrable final visé est
documenté dans \texttt{docs/PROTOCOLE\_EVALUATION.md}.

%==============================================================================
\section{Validation de la qualité du contenu généré}
\label{sec:val-contenu}

\subsection{Corpus produit et audit automatique}

La chaîne complète (perception simulée, décision, LLM + RAG) a produit
\textbf{220} contenus (\texttt{corpus\_200.csv}). Les métriques automatiques
donnent :

\begin{itemize}[leftmargin=1.4em]
  \item Longueur moyenne : \textbf{225} mots (min 115, max 434).
  \item Indice Kandel--Moles moyen : 94{,}2 (textes accessibles).
  \item Taux de répétition moyen : 0{,}485.
  \item Conformité à l'audit heuristique : \textbf{97{,}7\,\%} (215/220 sans
  anomalie majeure ; 5 contenus avec alertes mineures).
\end{itemize}

Les patrons d'encouragement (D1--D5) produisent les textes les plus courts
($\approx$130--190 mots) ; les exercices (C1--C2) les plus longs
($\approx$320--350 mots), en cohérence avec les contraintes des patrons.

\subsection{Évaluation experte (protocole)}

Un échantillon de 60 contenus stratifiés par type (A/B/C/D) a été extrait
(\texttt{corpus\_evaluation\_60.csv}). La grille comprend six critères Likert,
dont l'absence d'erreurs factuelles \emph{avec RAG}. Cette évaluation humaine
n'a pas été menée à terme dans le délai du mémoire ; elle constitue la
perspective immédiate de validation du fond pédagogique.

\subsection{Métriques automatiques complémentaires}

Le script \texttt{scripts/metrics\_contenu.py} calcule longueur, lisibilité et
présence heuristique des éléments structurants des patrons. Ces métriques
contrôlent la \emph{forme} mais ne garantissent pas la pertinence sémantique.

%==============================================================================
\section{Validation utilisateur (étude pilote)}
\label{sec:val-utilisateur}

\subsection{Protocole expérimental}

Un protocole \emph{within-subject} ($n=30$, deux conditions : contenu statique
vs contenu adaptatif \emolearn{}) est documenté (durée $\approx$45 min, SUS,
Likert utilité/adéquation émotionnelle). Les hypothèses exploratoires (C2 $>$
C1 sur utilité ; SUS $\geq 68$) n'ont pas été testées empiriquement dans cette
version du travail.

\subsection{Simulation de parcours adaptatifs}

Le prototype React permet néanmoins une démonstration qualitative : lecture de
séquences vidéo, affichage de l'état détecté, de la décision et du contenu
généré (avec sources RAG), export PDF et journal SQLite.

%==============================================================================
\section{Résultats et analyses}
\label{sec:val-resultats}

\subsection{Synthèse quantitative}

\begin{table}[H]
\centering
\small
\renewcommand{\arraystretch}{1.15}
\begin{tabular}{|>{\raggedright\arraybackslash}p{4.2cm}|c|c|}
\hline
\rowcolor{gray!20}
\textbf{Composant} & \textbf{Métrique clé} & \textbf{Valeur} \\
\hline
FER\,$\rightarrow$\,DAiSEE (appris) & Accuracy test & 0{,}723 \\
\hline
SER wav2vec2 (banc aligné) & Accuracy & 0{,}487 \\
\hline
Fusion vote pondéré & Accuracy & 0{,}500 \\
\hline
Fusion priorité pédagogique & Rappel urgents & 0{,}400 \\
\hline
Corpus généré (LLM + RAG) & Audit OK & 97{,}7\,\% \\
\hline
\end{tabular}
\caption{Synthèse des résultats chiffrés disponibles.}
\label{tab:synthese-val}
\end{table}

\subsection{Retours qualitatifs}

Les démonstrations internes du prototype mettent en évidence la \emph{lisibilité}
de la chaîne de décision (journal complet) et l'apport visible du RAG (sources
documentaires affichées dans l'interface). Les limites perceptives sur les
classes minoritaires et le compromis fusion accuracy/rappel urgent sont
explicitement documentées plutôt que masquées.

%==============================================================================
\section{Discussion approfondie}
\label{sec:val-discussion}

\subsection{Interprétation des résultats}

Les validations techniques confirment que (i) l'apprentissage léger améliore
fortement le mapping FER\,$\rightarrow$\,DAiSEE, (ii) le SER appris surpasse
l'heuristique vocale, (iii) la fusion apporte un gain mesurable mais impose un
arbitrage entre exactitude globale et sensibilité aux états critiques. Le module
RAG et les patrons structurés maintiennent la forme du corpus à un niveau
acceptable (97{,}7\,\% d'audit OK), sans préjuger de la qualité pédagogique
profonde, qui requiert l'évaluation experte prévue.

\subsection{Comparaison avec l'état de l'art}

Peu de travaux publient une chaîne comparable (multimodal + règles + RAG + LLM)
avec métriques de fusion \emph{et} corpus généré audité. La contribution
distinctive réside dans l'architecture reproductible et l'honnêteté sur les
compromis, plus que dans des scores de référence établis sur MOOC réels.

\subsection{Forces et limites identifiées}

\begin{itemize}[leftmargin=1.4em]
  \item \textbf{Forces} : pipeline bout-en-bout, traçabilité, RAG intégré,
  banc d'essai aligné, protocoles d'évaluation documentés.
  \item \textbf{Limites} : simulation (pas de captation live), appariement
  visage--voix hétérogène en production, évaluations experte et utilisateur non
  finalisées, déséquilibre \textsc{DAiSEE} pénalisant le F1 macro.
\end{itemize}
